% !PRINTED 255
\weberpage{255}

\begin{center}
\S.\,80.

\gesperrt{Die Trägheit der Formen zweiten Grades.}
\end{center}

Durch den Satz des vorigen Paragraphen sind wir auf die Untersuchung der homogenen Functionen zweiten Grades mit reellen Coëfficienten hingewiesen. Ueber diese Functionen haben wir in \S.\,58 unter dem Namen des Trägheitsgesetzes einen Satz kennen gelernt, der für das Folgende die Grundlage bildet.

Der Satz bestand darin, dass, wie man auch eine quadratische Form von $m$ Veränderlichen in eine Summe von positiven und negativen Quadraten von linear unabhängigen linearen Functionen transformiren mag, was auf unendlich viele verschiedene Arten möglich ist, die Anzahl der positiven und ebenso die der negativen Quadrate immer dieselbe ist.

Bezeichnen wir mit $\pi$ die Anzahl der positiven und mit $\nu$ die Anzahl der negativen Quadrate, so ist $\pi + \nu$ höchstens gleich $m$.

Wenn wir den Unterschied $m - \pi - \nu$ mit $\varrho$ bezeichnen, so kann, wenn wir die allgemeine quadratische Form von $m$ Variablen als Summe von $m$ Quadraten darstellen, $\varrho$ als die Anzahl der Quadrate mit verschwindenden Coëfficienten bezeichnet werden.

Die drei Zahlen $\pi$, $\nu$, $\varrho$, von denen keine negativ sein kann und deren Summe gleich der Anzahl der Variablen ist,
\weq{1}{\pi + \nu + \varrho = m,}
sind also für eine bestimmte quadratische Form unveränderlich.

Wir haben nach Mitteln zu suchen, um aus den Coëfficienten der quadratischen Form die Zahlen $\pi$, $\nu$, $\varrho$ zu ermitteln. Wenden wir diese Mittel auf die Bezoutiante an, so erhalten wir Kennzeichen, um aus den Coëfficienten einer Gleichung auf die Anzahl ihrer reellen Wurzeln zu schliessen.

Es sei also jetzt, wie in \S.\,56
\weq{2}{\varphi(x_1, x_2 \ldots x_m) = \sum_{1,m}^{i,k} a_{i,k} x_i x_k, \qquad a_{i,k} = a_{k,i}}
eine quadratische Form von $m$ Veränderlichen mit \gesperrt{reellen} Coëfficienten $a_{i,k}$.
